\documentclass{dayMath}
\begin{document}
\begin{problem}{2026.05.18}
计算下列求和.
\begin{enumerate}
\item \(S_n=1+2x+3x^2+\cdots+nx^{n-1},x\ne 0,n\in\mathbb{N}^+\).
\item \(S_n=C_n^1+2C_n^2+3C_n^3+\cdots+nC_n^n\).
\end{enumerate}
\end{problem}


\begin{problem}{2026.05.17}
对于\(\triangle ABC\)以及平面上一点\(O\), 有\(3\overrightarrow{OA}+4\overrightarrow{OB}+5\overrightarrow{OC}=\overrightarrow{0}\), 求\(\triangle OBC,\triangle OAC\)与\(\triangle OAB\)的面积之比.
\begin{solution}
\(3:4:5\).
\end{solution}
\end{problem}


\begin{problem}{2026.05.15}
过定点\(A(a,b)\)任作相互垂直的两条直线\(l_1\)与\(l_2\), 且\(l_1\)与\(x\)轴交于点\(M\), \(l_2\)与\(y\)轴交于点\(N\). 求\(MN\)中点\(P\)的轨迹方程.
\begin{solution}
\(2ax+2by-a^2-b^2=0\).
\end{solution}
\end{problem}


\begin{problem}{2026.05.12}
已知\(a,b,c\in\mathbb{R}^+\), 证明:
\[\frac{27a^2b^2c^2}{(a+b+c)^3}\geq (a+b-c)(a+c-b)(b+c-a).\]
\begin{solution}
当\((a+b-c)(a+c-b)(b+c-a)\leq 0\)时, 原不等式显然成立.

当\((a+b-c)(a+c-b)(b+c-a)> 0\)时, 以\(a,b,c\)为边恰好可构成\(\triangle ABC\). 根据三角形面积的海伦公式\(S=\sqrt{p(p-a)(p-b)(p-c)}\), 其中\(p=(a+b+c)/2\), 整理得\[S^2=\frac{1}{16}(a+b+c)(a+b-c)(b+c-a)(a+c-b),\]
带入原不等式即证明\(\frac{27a^2b^2c^2}{(a+b+c)^2}\geq 16S^2\), 即证\(\frac{3\sqrt{3}abc}{a+b+c}\geq 4S\), 即\[\frac{3\sqrt{3}abc}{a+b+c}\geq4\times\frac{1}{2}ab\sin C,\]
由正弦定理代入即证明\(\sin A+\sin B+\sin C\leq\frac{3\sqrt{3}}{2}\).
\end{solution}
\end{problem}
\begin{problem}{2026.05.11}
对于\(\triangle ABC\), 证明\(\sin A+\sin B+\sin C\leq\frac{3\sqrt{3}}{2}\).
\begin{solution}
由和差化积公式，
$$ \sin A + \sin B = 2\cos\frac{C}{2}\cos\frac{A-B}{2} \le 2\cos\frac{C}{2} $$
故
$$ \sin A + \sin B + \sin C \le 2\cos\frac{C}{2} + 2\sin\frac{C}{2}\cos\frac{C}{2} = 2\cos\frac{C}{2}\left(1+\sin\frac{C}{2}\right) $$

令 $x = 1+\sin\frac{C}{2}$，因 $C \in (0, \pi)$，故 $x \in (1, 2]$。此时 $\cos^2\frac{C}{2} = 1-\sin^2\frac{C}{2} = (2-x)x$。

将上式平方得：
$$ \left[2\cos\frac{C}{2}\left(1+\sin\frac{C}{2}\right)\right]^2 = 4(2-x)x \cdot x^2 = 4(2-x)x^3 $$

由均值不等式，因 $x \in (1, 2]$，
$$ (2-x) + \frac{x}{3} + \frac{x}{3} + \frac{x}{3} = 2 \ge 4\sqrt[4]{(2-x)\left(\frac{x}{3}\right)^3} $$
解得 $(2-x)x^3 \le \frac{27}{16}$，从而
$$ 4(2-x)x^3 \le 4 \times \frac{27}{16} = \frac{27}{4} $$

开方即得 $\sin A + \sin B + \sin C \le \frac{3\sqrt{3}}{2}$，当且仅当 $A=B=C=\frac{\pi}{3}$ 时等号成立。
\end{solution}
\end{problem}
\end{document}